\documentclass[12pt,letterpaper]{article}
\usepackage{graphicx,textcomp}
\usepackage{natbib}
\usepackage{setspace}
\usepackage{fullpage}
\usepackage{color}
\usepackage[reqno]{amsmath}
\usepackage{amsthm}
\usepackage{fancyvrb}
\usepackage{amssymb,enumerate}
\usepackage[all]{xy}
\usepackage{endnotes}
\usepackage{lscape}
\newtheorem{com}{Comment}
\usepackage{float}
\usepackage{hyperref}
\newtheorem{lem} {Lemma}
\newtheorem{prop}{Proposition}
\newtheorem{thm}{Theorem}
\newtheorem{defn}{Definition}
\newtheorem{cor}{Corollary}
\newtheorem{obs}{Observation}
\usepackage[compact]{titlesec}
\usepackage{dcolumn}
\usepackage{tikz}
\usetikzlibrary{arrows}
\usepackage{multirow}
\usepackage{xcolor}
\newcolumntype{.}{D{.}{.}{-1}}
\newcolumntype{d}[1]{D{.}{.}{#1}}
\definecolor{light-gray}{gray}{0.65}
\usepackage{url}
\usepackage{listings}
\usepackage{color}

\definecolor{codegreen}{rgb}{0,0.6,0}
\definecolor{codegray}{rgb}{0.5,0.5,0.5}
\definecolor{codepurple}{rgb}{0.58,0,0.82}
\definecolor{backcolour}{rgb}{0.95,0.95,0.92}

\lstdefinestyle{mystyle}{
	backgroundcolor=\color{backcolour},   
	commentstyle=\color{codegreen},
	keywordstyle=\color{magenta},
	numberstyle=\tiny\color{codegray},
	stringstyle=\color{codepurple},
	basicstyle=\footnotesize,
	breakatwhitespace=false,         
	breaklines=true,                 
	captionpos=b,                    
	keepspaces=true,                 
	numbers=left,                    
	numbersep=5pt,                  
	showspaces=false,                
	showstringspaces=false,
	showtabs=false,                  
	tabsize=2
}
\lstset{style=mystyle}
\newcommand{\Sref}[1]{Section~\ref{#1}}
\newtheorem{hyp}{Hypothesis}

\title{Problem Set 3 - Applied Stats II}
\date{Due: March 25, 2026}
\author{Sarah Magdihs}


\begin{document}
	\maketitle
	\section*{Instructions}
	\begin{itemize}
	\item Please show your work! You may lose points by simply writing in the answer. If the problem requires you to execute commands in \texttt{R}, please include the code you used to get your answers. Please also include the \texttt{.R} file that contains your code. If you are not sure if work needs to be shown for a particular problem, please ask.
\item Your homework should be submitted electronically on GitHub in \texttt{.pdf} form.
\item This problem set is due before 23:59 on Wednesday March 25, 2026. No late assignments will be accepted.
	\end{itemize}

	\vspace{.25cm}
\section*{Question 1}
\vspace{.25cm}
\noindent We are interested in how governments' management of public resources impacts economic prosperity. Our data come from \href{https://www.researchgate.net/profile/Adam_Przeworski/publication/240357392_Classifying_Political_Regimes/links/0deec532194849aefa000000/Classifying-Political-Regimes.pdf}{Alvarez, Cheibub, Limongi, and Przeworski (1996)} and is labelled \texttt{gdpChange.csv} on GitHub. The dataset covers 135 countries observed between 1950 or the year of independence or the first year forwhich data on economic growth are available ("entry year"), and 1990 or the last year for which data on economic growth are available ("exit year"). The unit of analysis is a particular country during a particular year, for a total $>$ 3,500 observations. 

\begin{itemize}
	\item
	Response variable: 
	\begin{itemize}
		\item \texttt{GDPWdiff}: Difference in GDP between year $t$ and $t-1$. Possible categories include: "positive", "negative", or "no change"
	\end{itemize}
	\item
	Explanatory variables: 
	\begin{itemize}
		\item
		\texttt{REG}: 1=Democracy; 0=Non-Democracy
		\item
		\texttt{OIL}: 1=if the average ratio of fuel exports to total exports in 1984-86 exceeded 50\%; 0= otherwise
	\end{itemize}
	
\end{itemize}
\newpage
\noindent Please answer the following questions:

\begin{enumerate}
	\item Construct and interpret an unordered multinomial logit with \texttt{GDPWdiff} as the output and "no change" as the reference category, including the estimated cutoff points and coefficients.
	\item Construct and interpret an ordered multinomial logit with \texttt{GDPWdiff} as the outcome variable, including the estimated cutoff points and coefficients.
	
	
\end{enumerate}

\section*{Answers to Question 1}
\vspace{.25cm}

\noindent\textbf{Task 1: Unordered Multinomial Logit}\\
\noindent To begin, I took a closer look at the dataset and recoded the relevant variables. Specifically, I made sure that  \texttt{GDPWdiff}  would be a factor with three levels. I set the reference category to be \texttt{"no change"}, as required by the task.  Additionally, I recoded the two explanatory variables as factors and gave them new labels. 

\noindent\textit{R Code:}
\lstinputlisting[language=R, firstline=50, lastline=75]{PS3_answersSM.R}  

\noindent Then, I specified the regression model. The output is presented in Table 1.\\
\noindent\textit{R Code:}
\lstinputlisting[language=R, firstline=79, lastline=86]{PS3_answersSM.R}  
\lstinputlisting[language=R, firstline=92, lastline=93]{PS3_answersSM.R}  
% Table created by stargazer v.5.2.3 by Marek Hlavac, Social Policy Institute. E-mail: marek.hlavac at gmail.com
% Date and time: Wed, Mar 25, 2026 - 15:02:26
\begin{table}[H] \centering 
	\caption{Unordered Multinomial Logit} 
	\label{} 
	\begin{tabular}{@{\extracolsep{5pt}}lcc} 
		\\[-1.8ex]\hline 
		\hline \\[-1.8ex] 
		& \multicolumn{2}{c}{\textit{Dependent variable:}} \\ 
		\cline{2-3} 
		\\[-1.8ex] & negative  & positive \\ 
		\\[-1.8ex] & (vs. no change) & (vs. no change) \\ 
		\\[-1.8ex] & (1) & (2)\\ 
		\hline \\[-1.8ex] 
		REG (Democracy) & 1.379$^{*}$ & 1.769$^{**}$ \\ 
		& (0.769) & (0.767) \\ 
		& & \\ 
		OIL (Fuel exports over 50\%) & 4.784 & 4.576 \\ 
		& (6.885) & (6.885) \\ 
		& & \\ 
		Constant (Intercept) & 3.805$^{***}$ & 4.534$^{***}$ \\ 
		& (0.271) & (0.269) \\ 
		& & \\ 
		\hline \\[-1.8ex] 
		Akaike Inf. Crit. & 4,690.770 & 4,690.770 \\ 
		\hline 
		\hline \\[-1.8ex] 
		\textit{Note:}  & \multicolumn{2}{r}{$^{*}$p$<$0.1; $^{**}$p$<$0.05; $^{***}$p$<$0.01} \\ 
	\end{tabular} 
\end{table}

% Table created by stargazer v.5.2.3 by Marek Hlavac, Social Policy Institute. E-mail: marek.hlavac at gmail.com
% Date and time: Wed, Mar 25, 2026 - 15:28:09
\begin{table}[H] \centering 
	\caption{OR for the Unordered Multinomial Logit} 
	\label{} 
	\begin{tabular}{@{\extracolsep{5pt}} cccc} 
		\\[-1.8ex]\hline 
		\hline \\[-1.8ex] 
		& (Intercept) & REG (Democracy) & OIL (Fuel exports over 50\%) \\ 
		\hline \\[-1.8ex] 
		negative & $44.942$ & $3.972$ & $119.578$ \\ 
		positive & $93.108$ & $5.865$ & $97.156$ \\ 
		\hline \\[-1.8ex] 
	\end{tabular} 
\end{table}

\noindent Fundamentally, the unordered multinomial logit regression models the log-odds of
being in either the “negative” or “positive” category relative to the baseline category (“no change”).

\noindent \textbf{Interpretation:}

\noindent\textit{\textbf{Constant:}}
\noindent For both models, the results show that the Intercept is positive and statistically significant. As further shown in Table 2, this means that even when all explanatory variables equal zero (meaning: Non-Democracy and the average ratio of fuel exports to total exports in 1984-86 did not exceed 50\% ), experiencing either a positive or negative change is much more likely than there being no change. Notably, the baseline odds are quite high. However, this could be the case because there are actually very few cases of "no change" in comparison to "positive/negative change".

\noindent\textit{\textbf{REG:}}\\
\noindent \texttt{negative vs. no change:} The coefficient for "REG" is $1.379$ (p $<$ 0.1). This means that, holding all other variables constant, being a democracy (rather than a non-democracy) is associated with a $1.379$ increase in the log-odds of experiencing a negative GDP change (relative to no change). Put differently, being a democracy (rather than a non-democracy) multiplies the odds of experiencing a negative GDP change (relative to no change) by a factor of exp($1.379$); meaning the odds are roughly $3.97$ times higher.

\noindent \texttt{positive vs. no change:} The coefficient for "REG" is $1.769$ (p $<$ 0.05). This means that, holding all other variables constant, being a democracy (rather than a non-democracy) is associated with a $1.769$ increase in the log-odds of experiencing a positive GDP change (relative to no change). Put differently, democracies are about $5.87$ times more likely to experience a positive change in GDP (relative to no change) than non-democracies.\\
\noindent\textit{\textbf{OIL:}} Based on the output for both comparisons, "OIL" does not have a statistically significant effect and there is \textit{a lot} of uncertainty concerning the coefficients. (Technically, the interpretation would follow the same 'structure' as the interpretation for REG).\\


\noindent\textbf{Task 2: Ordered Multinomial Logit}
\vspace{.25cm}
\\As outlined by the assignment, I also constructed an ordered multinomial logit regression. To do so, I first made sure that the response variable would be an ordered factor. Then, I ran the regression. The output is shown in the Tables below. \\
\noindent\textit{R Code:}
\lstinputlisting[language=R, firstline=98, lastline=127]{PS3_answersSM.R}  


% Table created by stargazer v.5.2.3 by Marek Hlavac, Social Policy Institute. E-mail: marek.hlavac at gmail.com
% Date and time: Wed, Mar 25, 2026 - 16:09:36
\begin{table}[H] \centering 
	\caption{Ordered Multinomial Logit} 
	\label{} 
	\begin{tabular}{@{\extracolsep{5pt}}lc} 
		\\[-1.8ex]\hline 
		\hline \\[-1.8ex] 
		& \multicolumn{1}{c}{\textit{Dependent variable:}} \\ 
		\cline{2-2} 
		\\[-1.8ex] & Difference in GDP (GDPWdiff)\\ 
		\\[-1.8ex] & negative change $<$ no change $<$ positive change \\ 
		\hline \\[-1.8ex] 
		REG (Democracy) & 0.398$^{***}$ \\ 
		& (0.075) \\ 
		& \\ 
		OIL (Fuel exports over 50\%) & $-$0.199$^{*}$ \\ 
		& (0.116) \\ 
		& \\ 
		\hline \\[-1.8ex] 
		Observations & 3,721 \\ 
		\hline 
		\hline \\[-1.8ex] 
		\textit{Note:}  & \multicolumn{1}{r}{$^{*}$p$<$0.1; $^{**}$p$<$0.05; $^{***}$p$<$0.01} \\ 
	\end{tabular} 
\end{table} 

% Table created by stargazer v.5.2.3 by Marek Hlavac, Social Policy Institute. E-mail: marek.hlavac at gmail.com
% Date and time: Wed, Mar 25, 2026 - 16:43:02
\begin{table}[H] \centering 
	\caption{Cutoffs} 
	\label{} 
	\begin{tabular}{@{\extracolsep{5pt}} cc} 
		\\[-1.8ex]\hline 
		\hline \\[-1.8ex] 
		negative\textbar no change & no change\textbar positive \\ 
		\hline \\[-1.8ex] 
		$$-$0.731$ & $$-$0.710$ \\ 
		\hline \\[-1.8ex] 
	\end{tabular} 
\end{table}

% Table created by stargazer v.5.2.3 by Marek Hlavac, Social Policy Institute. E-mail: marek.hlavac at gmail.com
% Date and time: Wed, Mar 25, 2026 - 16:09:36
\begin{table}[H] \centering 
	\caption{OR: for the Ordered Multinomial Logit} 
	\label{} 
	\begin{tabular}{@{\extracolsep{5pt}} cccc} 
		\\[-1.8ex]\hline 
		\hline \\[-1.8ex] 
		& OR & 2.5 \% & 97.5 \% \\ 
		\hline \\[-1.8ex] 
		REG (Democracy)  & $1.490$ & $1.286$ & $1.727$ \\ 
		OIL (Fuel exports over 50\%) & $0.820$ & $0.655$ & $1.031$ \\ 
		\hline \\[-1.8ex] 
	\end{tabular} 
\end{table}

\noindent Whereas the unordered model compared each category to the reference category, the ordered model follows a cumulative logic . As shown in Table 3, the variable is ordered as \texttt{negative change $<$ no change $<$ positive change}. Crucially, the ordered model assumes that the relationship between each pair of outcome groups is equal (proportional odds). Hence, there is only one set of coefficients\\
\noindent \textbf{Interpretation:}\\
\noindent\textit{\textbf{Cutoff Points:}}
\noindent The first cutoff point (\texttt{negative|no change}) is $-0.731$, whereas the second cutoff point (\texttt{no change|positive}) is $-0.710$. Notably, these two points are rather close to one another, which could be a reflection of earlier results that experiencing either a positive or negative change is much more likely than there being no change.\\
\noindent\textit{\textbf{REG:}}
\noindent The coefficient for REG is $0.398$ and is statistically significant (p $<$ 0.01). Thus, holding other variables constant, democracies are about $1.49$ times more likely to be in a higher category of GDP change (no change/positive change). Given the proportional odds assumption, this effect applies to both thresholds.\\
\noindent\textit{\textbf{OIL:}}
\noindent Unlike in the unordered model, the coefficient for OIL ($-0.199$) is statistically significant in the ordered model (p $<$ 0.1). Thus, holding other variables constant, countries where the average ratio of fuel exports to total exports in 1984-86 exceeded 50\% have about 18\% lower odds of being in a higher category of GDP change (no change/positive change).

\section*{Question 2} 
\vspace{.25cm}

\noindent Consider the data set \texttt{MexicoMuniData.csv}, which includes municipal-level information from Mexico. The outcome of interest is the number of times the winning PAN presidential candidate in 2006 (\texttt{PAN.visits.06}) visited a district leading up to the 2009 federal elections, which is a count. Our main predictor of interest is whether the district was highly contested, or whether it was not (the PAN or their opponents have electoral security) in the previous federal elections during 2000 (\texttt{competitive.district}), which is binary (1=close/swing district, 0="safe seat"). We also include \texttt{marginality.06} (a measure of poverty) and \texttt{PAN.governor.06} (a dummy for whether the state has a PAN-affiliated governor) as additional control variables. 

\begin{enumerate}
	\item [(a)]
	Run a Poisson regression because the outcome is a count variable. Is there evidence that PAN presidential candidates visit swing districts more? Provide a test statistic and p-value.

	\item [(b)]
	Interpret the \texttt{marginality.06} and \texttt{PAN.governor.06} coefficients.
	
	\item [(c)]
	Provide the estimated mean number of visits from the winning PAN presidential candidate for a hypothetical district that was competitive (\texttt{competitive.district}=1), had an average poverty level (\texttt{marginality.06} = 0), and a PAN governor (\texttt{PAN.governor.06}=1).
	
\end{enumerate}


\section*{Answers to Question 2}
\vspace{.25cm}
\noindent\textbf{Task a: Poisson Regression}\\
First, I took a quick look at the dataset and familiarised myself with the variables.Then, I specified the regression model and calculated the test statistic and p value. The outputs are presented below the R Code.\\
\textit{R Code:}
\lstinputlisting[language=R, firstline=133, lastline=140]{PS3_answersSM.R}  
\lstinputlisting[language=R, firstline=142, lastline=155]{PS3_answersSM.R}  

% Table created by stargazer v.5.2.3 by Marek Hlavac, Social Policy Institute. E-mail: marek.hlavac at gmail.com
% Date and time: Wed, Mar 25, 2026 - 17:42:36
\begin{table}[H] \centering 
	\caption{Poisson Regression Model} 
	\label{} 
	\begin{tabular}{@{\extracolsep{5pt}}lc} 
		\\[-1.8ex]\hline 
		\hline \\[-1.8ex] 
		& \multicolumn{1}{c}{\textit{Dependent variable:}} \\ 
		\cline{2-2} 
		\\[-1.8ex] & PAN.visits.06 \\ 
		\hline \\[-1.8ex] 
		competitive.district & $-$0.081 \\ 
		& (0.171) \\ 
		& \\ 
		marginality.06 & $-$2.080$^{***}$ \\ 
		& (0.117) \\ 
		& \\ 
		PAN.governor.06 & $-$0.312$^{*}$ \\ 
		& (0.167) \\ 
		& \\ 
		Constant & $-$3.810$^{***}$ \\ 
		& (0.222) \\ 
		& \\ 
		\hline \\[-1.8ex] 
		Observations & 2,407 \\ 
		Log Likelihood & $-$645.606 \\ 
		Akaike Inf. Crit. & 1,299.213 \\ 
		\hline 
		\hline \\[-1.8ex] 
		\textit{Note:}  & \multicolumn{1}{r}{$^{*}$p$<$0.1; $^{**}$p$<$0.05; $^{***}$p$<$0.01} \\ 
	\end{tabular} 
\end{table} 

% Table created by stargazer v.5.2.3 by Marek Hlavac, Social Policy Institute. E-mail: marek.hlavac at gmail.com
% Date and time: Wed, Mar 25, 2026 - 17:44:16
\begin{table}[H] \centering 
	\caption{Test statistic and p value for competitive.districts} 
	\label{} 
	\begin{tabular}{@{\extracolsep{5pt}} cccc} 
		\\[-1.8ex]\hline 
		\hline \\[-1.8ex] 
		Estimate & Std. Error & z value & Pr(\textgreater \textbar z\textbar ) \\ 
		\hline \\[-1.8ex] 
		$$-$0.081$ & $0.171$ & $$-$0.477$ & $0.634$ \\ 
		\hline \\[-1.8ex] 
	\end{tabular} 
\end{table} 

\noindent The coefficient for \texttt{competitive.district} is $-0.081$ with a z-statistic of $-0.477$ and a p-value of $0.634$. Therefore, the model shows no statistically significant evidence that PAN presidential candidates visited swing districts more. Although the coefficient is negative, the estimated effect is not distinguishable from zero.


\noindent\textbf{Task b: Interpretation}\\
To make the interpretation easier, I exponentiated the coefficients. \\
\textit{R Code:}
\lstinputlisting[language=R, firstline=156, lastline=158]{PS3_answersSM.R}  
% Table created by stargazer v.5.2.3 by Marek Hlavac, Social Policy Institute. E-mail: marek.hlavac at gmail.com
% Date and time: Wed, Mar 25, 2026 - 18:01:04
\begin{table}[H] \centering 
	\caption{} 
	\label{} 
	\begin{tabular}{@{\extracolsep{5pt}} cccc} 
		\\[-1.8ex]\hline 
		\hline \\[-1.8ex] 
		(Intercept) & competitive.district & marginality.06 & PAN.governor.06 \\ 
		\hline \\[-1.8ex] 
		$0.022$ & $0.922$ & $0.125$ & $0.732$ \\ 
		\hline \\[-1.8ex] 
	\end{tabular} 
\end{table} 
\noindent\textit{\textbf{marginality.06}:} The coefficient is $-2.080$ and is statistically significant (p $<$ 0.01). This means that, holding all other variables constant, a one-unit increase in marginality.06 (poverty) is associated with a decrease of $2.080$ in the log expected number of visits. Put differently, a one-unit increase in marginality.06 reduces the expected number of visits by 87.5\%.\\
\noindent\textit{\textbf{PAN.governor.06}:} The coefficient is $-0.312$ and is statistically significant (p $<$ 0.1). This means that, holding all other variables constant, having a PAN-affiliated governor is associated with a decrease of $-0.312$ in the log expected number of visits. Put differently, having a PAN-affiliated governor reduces the expected number of visits by 26.8\%.\\


\noindent\textbf{Task c: Estimated mean number of visits for a specific scenario}\\
\noindent Lastly, I estimated the mean number of visits for a hypothetical district that was competitive, had an average poverty level, and a PAN governor. As shown below, the estimated mean number of visits for the hypothetical district is approximately $0.015$. \\
\textit{R Code:}
\lstinputlisting[language=R, firstline=164, lastline=170]{PS3_answersSM.R}  

\begin{verbatim}
> predict(poisson_reg, newdata = predict_case, type = "response")
1 
0.01494818 
\end{verbatim} 


\end{document}
